\documentclass[11pt,twocolumn]{article}

%% Reasoning-aware Adaptive Routing for Hybrid GraphRAG
%% Nguyen Dinh Thanh - MSSV: 20225670
%% School of Information Technology and Communication
%% Hanoi University of Science and Technology
%% Supervisors: Pham Van Hai & Nguyen Viet Cuong

% ============================================================
% Encoding and font
% ============================================================
\usepackage[utf8]{inputenc}
\usepackage[T5]{fontenc}
\usepackage[vietnamese,english]{babel}
\usepackage{lmodern}

% ============================================================
% Math, tables, algorithms, figures
% ============================================================
\usepackage{graphicx}
\usepackage{amssymb}
\usepackage{amsmath}
\usepackage{amsthm}
\usepackage{mathtools}
\usepackage{mathrsfs}
\usepackage{scrextend}
\usepackage{multirow}
\usepackage{xcolor}
\usepackage{booktabs}
\usepackage{tabularx}
\usepackage{array}
\usepackage{ragged2e}
\usepackage{makecell}
\usepackage{adjustbox}
\usepackage{algorithm}
\usepackage{algorithmicx}
\usepackage{algpseudocode}
\usepackage{float}
\usepackage{placeins}

% ============================================================
% TikZ
% ============================================================
\usepackage{tikz}
\usetikzlibrary{arrows.meta, positioning, shapes.geometric, fit, calc}

% ============================================================
% Layout control for two-column paper
% ============================================================
\usepackage{titlesec}
\raggedbottom
\setlength{\parskip}{0pt}
\setlength{\parindent}{1em}

\titlespacing*{\section}
  {0pt}{1.0ex plus 0.3ex minus 0.2ex}{0.7ex plus 0.2ex}

\titlespacing*{\subsection}
  {0pt}{0.8ex plus 0.25ex minus 0.2ex}{0.45ex plus 0.15ex}

\titlespacing*{\subsubsection}
  {0pt}{0.7ex plus 0.2ex minus 0.2ex}{0.35ex plus 0.1ex}

% ============================================================
% Hyperref should be loaded near the end
% ============================================================
%\usepackage{microtype}
\usepackage{hyperref}

\usepackage{pgfplots}
\pgfplotsset{compat=1.18}

% ============================================================
% Custom column types
% ============================================================
\newcolumntype{Y}{>{\RaggedRight\arraybackslash}X}
\newcolumntype{L}[1]{>{\RaggedRight\arraybackslash}p{#1}}

% ============================================================
% Misc.
% ============================================================
\hbadness=10000
\vbadness=10000
\newtheorem{proposition}{Proposition}

\def\abstract{
\typeout{Abstract}
 {\bf Abstract}
} \def\endabstract{\par}

\begin{document}

\title{Reasoning-aware Adaptive Routing for Hybrid GraphRAG}

\author{
  Nguyen Dinh Thanh \\
  School of Information Technology and Communication, \\
  Hanoi University of Science and Technology, \\
  Hai Ba Trung, Hanoi, Vietnam \\
  \texttt{thanh.nd225670@sis.hust.edu.vn}
}

\maketitle

% ============================================================
% Abstract
% ============================================================

\begin{abstract}
Retrieval-Augmented Generation improves the factual
reliability of large language models by grounding generated
answers in external knowledge sources. In Vietnamese legal
question answering, however, a fixed retrieval strategy is
often insufficient. Some questions can be answered through
direct textual lookup, while others require reasoning over
legal relations such as amendment, repeal, replacement,
cross-document reference, effective date, issuing authority,
and article membership. Moreover, when a query is ambiguous
or underspecified, immediate retrieval may ground the answer
in the wrong legal document, article, entity, or time period.

This paper proposes a Vietnamese legal question answering
system based on Hybrid GraphRAG with a two-stage
Reasoning-aware Adaptive Router. Given a user query and an
optional conversation history, the router selects one of four
actions: dense retrieval, graph traversal, hybrid reasoning,
or clarification. Stage~1 uses an XGBoost classifier over
query-level reasoning and ambiguity features to make fast
routing decisions. Stage~2 invokes an LLM-based verifier only
for uncertain, ambiguous, or relation-heavy cases. The
retrieval backend combines a dense vector index over
Vietnamese legal text chunks with a Neo4j legal knowledge
graph containing legal documents, legal articles, vector
chunks, agencies, concepts, and legal relations.

Experiments on a PD-only strict $600$-query benchmark
(\texttt{phapdien\_strict}) show that routing is
template-separable: the Single-stage Router and the
Two-stage Hybrid both reach routing accuracy $1.0000$.
However, end-to-end answer F1 is dominated by dense retrieval
performance: Pure Vector achieves F1 $0.4947$, while Pure Graph
achieves F1 $0.3602$. The deployed router systems remain close
to the graph-only baseline (Single-stage Router F1 $0.3610$,
Two-stage Hybrid F1 $0.3602$), because graph/hybrid routes
often retrieve insufficient evidence for the canonical targets.
On the original constructed ambiguity benchmark, the initial
Stage~2 configuration increases clarification F1 to $0.684$
with precision $1.000$. After adding deterministic
history-referent resolution and a severe-ambiguity override
policy, the same ambiguity benchmark improves to
clarification F1 $=0.840$. On a separate
$160$-query conversation stress test, the router reaches
clarification F1 $=0.862$ and history-resolution accuracy
$=0.850$. These diagnostic results do not replace the strict
benchmark; they show that ambiguity handling improves while
multi-interpretation cases and false clarification calibration
remain open issues.

\textbf{Keywords:} Retrieval-Augmented Generation,
GraphRAG, Hybrid RAG, Adaptive Routing, Vietnamese Legal QA,
Knowledge Graphs, Multi-hop Reasoning, Ambiguity Detection.
\end{abstract}

% ============================================================
\section{Introduction}
\label{sec:intro}
% ============================================================

Large language models (LLMs) have demonstrated strong
language understanding and generation ability, but their use
in knowledge-intensive question answering remains limited
when answers must be grounded in external, updateable, and
verifiable evidence rather than relying only on model
parameters \cite{lewis2020rag,guu2020realm}. Retrieval-
Augmented Generation (RAG) addresses this limitation by
retrieving relevant passages from an external corpus and
conditioning the generator on the retrieved evidence, which
makes the system more suitable for domains where factual
knowledge changes over time \cite{karpukhin2020dpr,
asai2023selfrag}. However, conventional dense-retrieval RAG
represents a corpus mainly as independent text chunks and is
therefore more effective for single-hop or direct lookup
questions than for questions requiring explicit reasoning
over relations among documents, entities, or conditions
\cite{trivedi2023interleaving,press2023measuring}. This
limitation is especially important in legal question
answering, where a user query may require identifying a
specific article, following a cross-document reference,
checking whether a regulation has been amended or repealed,
or reasoning over issuing authorities and effective dates
\cite{chalkidis2020legalbert,zhong2020nlp_law,
sovrano2020legalkg}. In Vietnamese legal QA, this problem
is further complicated by Vietnamese-specific preprocessing
issues such as word segmentation and multi-syllable named
entities, which affect entity detection and legal reference
matching \cite{nguyen2020phobert,vu2018vncorenlp}.

Recent work has explored several directions to improve RAG.
Self-RAG introduces reflection signals to decide when
retrieval is needed and whether retrieved passages are
useful, while FLARE performs active retrieval during
generation when missing information is anticipated
\cite{asai2023selfrag,jiang2023flare}. RAPTOR organizes
documents into hierarchical summaries for multi-level
retrieval, improving retrieval over long or complex
document collections \cite{sarthi2024raptor}. Graph-based
retrieval and GraphRAG further address relational retrieval
by representing documents, entities, and relations as graph
structures, allowing retrieval to follow explicit links
rather than relying only on vector similarity
\cite{edge2024graphrag,wang2024kgp,li2024graphreader}.
Adaptive-RAG shows that different questions may require
different retrieval and reasoning strategies depending on
question complexity, while FrugalGPT and RouteLLM show that
routing can reduce inference cost by assigning inputs to
suitable models or inference paths
\cite{jeong2024adaptive,chen2023frugalgpt,ong2024routellm}.
Ambiguity-aware QA research also shows that underspecified
or context-dependent questions may require rewriting or
clarification before reliable answering
\cite{min2020ambigqa,elgohary2019canard}. The difficulty is
that these actions have different cost and reasoning
profiles that cannot be reduced to a single routing
dimension: dense retrieval is fast but relationally blind,
graph traversal is more expressive but computationally
heavier, hybrid reasoning combines both evidence sources at
higher cost, and clarification introduces an additional
conversational turn before retrieval. Nevertheless, existing
directions do not directly address the four-way decision
needed in Vietnamese Legal Hybrid GraphRAG: whether a query
should use dense retrieval, graph traversal, hybrid
reasoning, or clarification before retrieval.

Motivated by this gap, this paper proposes a
Reasoning-aware Adaptive Router for Vietnamese Legal Hybrid
GraphRAG. The central challenge is that legal queries vary
in both reasoning depth and ambiguity: a short query may
still require multi-hop legal reasoning, while a long query
may only need direct article lookup; similarly, an
underspecified query can lead the system to retrieve the
wrong document, article, or time period if clarification is
not performed first. The motivation of this work is therefore
to make retrieval strategy selection an explicit intermediate
decision before answer generation. Given a user query and
optional conversation history, the proposed router first
extracts reasoning and ambiguity features, then uses a fast
Stage~1 XGBoost classifier to predict an initial route, and
finally invokes a Stage~2 LLM verifier only for uncertain,
ambiguous, or relation-heavy cases
\cite{chen2016xgboost,jeong2024adaptive}.

The main contributions of this paper are three-fold. First,
we formulate four-way reasoning-aware routing for Vietnamese
Legal Hybrid GraphRAG, where the system must choose among
dense retrieval, graph traversal, hybrid reasoning, and
clarification before answer generation. Second, we propose a
two-stage cost-aware router that combines a lightweight
feature-based Stage~1 classifier with a selective LLM-based
Stage~2 verifier for uncertain, ambiguous, or relation-heavy
queries. Third, we evaluate the system on a strict Vietnamese
legal QA benchmark and two diagnostic ambiguity benchmarks,
separating end-to-end retrieval-route selection from
clarification and conversation-history behavior. This
separation is important because a router may improve safe
clarification on ambiguous inputs without necessarily
improving strict routing accuracy on a benchmark that
contains no intended clarification queries.

The remainder of the paper is organized as follows.
Section~\ref{sec:related_work} reviews related work and
positions the proposed router. Section~\ref{sec:methodology}
defines the routing problem and describes the system design.
Sections~\ref{sec:eval} and~\ref{sec:results_discussion}
present the experimental setup and results. Finally,
Section~\ref{sec:limitations} discusses limitations and
future work before Section~\ref{sec:conclusion} concludes the
paper.

% ============================================================
\section{Related Work}
\label{sec:related_work}
% ============================================================

The proposed system is related to six research lines:
retrieval-augmented generation, graph retrieval and
GraphRAG, adaptive routing, ambiguity handling, dataset and
benchmark construction, and Vietnamese legal NLP.

\subsection{Retrieval-Augmented Generation}

Retrieval-Augmented Generation combines parametric
knowledge from neural language models with non-parametric
knowledge retrieved from an external corpus
\cite{lewis2020rag,guu2020realm}. In a standard RAG
pipeline, the input query is encoded into a dense
representation, relevant passages are retrieved from a
vector index, and the generator produces an answer
conditioned on the retrieved context \cite{karpukhin2020dpr}.
Dense passage retrieval is effective for open-domain QA
because it can retrieve semantically relevant passages even
when lexical overlap is limited. However, it usually treats
the corpus as a collection of independent chunks and is
therefore less suitable for questions requiring explicit
reasoning over document structure, entity relations, or
multi-hop evidence chains
\cite{trivedi2023interleaving,press2023measuring}.

Recent RAG variants attempt to improve retrieval control and
evidence usage. Self-RAG learns to retrieve, generate, and
critique through self-reflection \cite{asai2023selfrag};
FLARE actively retrieves additional evidence during
generation \cite{jiang2023flare}; and RAPTOR constructs
tree-organized summaries for retrieval at different
abstraction levels \cite{sarthi2024raptor}. These methods
improve factual grounding, but they do not explicitly
formulate the problem of choosing among dense retrieval,
graph traversal, hybrid reasoning, and clarification in a
legal Hybrid GraphRAG system.

\subsection{Graph Retrieval and GraphRAG}

Graph-based retrieval represents knowledge as nodes and
edges, enabling retrieval to follow explicit relationships
rather than relying only on vector similarity
\cite{sun2019pullnet,wang2024kgp}. Microsoft GraphRAG
builds an entity graph and community summaries for
query-focused summarization over private corpora
\cite{edge2024graphrag}. Knowledge Graph Prompting
constructs a graph over multiple documents and uses graph
traversal to gather supporting passages for multi-document
QA \cite{wang2024kgp}. GraphReader structures long documents
as a graph and allows an agent to explore relevant nodes and
neighbors \cite{li2024graphreader}. HippoRAG proposes a
neurobiologically inspired graph memory for RAG, using a
hippocampus-like indexing scheme to support multi-hop
retrieval over long-term knowledge
\cite{gutierrez2024hipporag}. While this approach improves
associative retrieval, it does not address the routing
decision of when graph traversal is necessary versus direct
dense lookup.

For legal QA, graph-based retrieval is particularly relevant
because legal documents naturally contain structured
relations such as amendment, repeal, replacement, reference,
jurisdiction, issuing authority, article membership, and
effective date \cite{sovrano2020legalkg}. However, graph
construction and traversal introduce additional
computational cost because the system must perform entity
linking, relation matching, graph querying, and evidence
serialization. Therefore, graph retrieval should be used
selectively rather than applied to every query.

\subsection{Adaptive Routing and Cost-aware Inference}

Adaptive routing aims to choose the most appropriate model,
retriever, or reasoning strategy for each input. Adaptive-RAG
proposes selecting QA strategies based on question
complexity, showing that not all questions require the same
retrieval or reasoning procedure \cite{jeong2024adaptive}.
FrugalGPT and RouteLLM demonstrate a related cost-aware
principle at the model level: when models or inference paths
have different costs and capabilities, routing can reduce
inference cost while maintaining quality
\cite{chen2023frugalgpt,ong2024routellm}. These works
motivate the cost-aware aspect of our router, but the routing
target in Vietnamese legal Hybrid GraphRAG is different.
Instead of only choosing a model size, retrieval depth, or
single QA strategy, the proposed system must decide which
evidence mechanism should be activated before answer
generation. Dense retrieval is suitable for direct textual
lookup, graph traversal is suitable for relation-oriented
legal questions, hybrid reasoning is useful when both article
text and graph paths are needed, and clarification is the
safer action when the legal target is underspecified. Thus,
our router is designed as a lightweight gateway that selects
among evidence mechanisms and clarification actions under
both reasoning and cost constraints.

\subsection{Ambiguity Handling in QA}

Real user questions are often ambiguous, incomplete, or
dependent on previous conversational turns. CANARD studies
question rewriting for context-dependent questions
\cite{elgohary2019canard}, while AmbigQA shows that
open-domain questions may have multiple valid
interpretations that should be handled explicitly
\cite{min2020ambigqa}. In high-stakes domains such as law,
ambiguity is especially problematic because answering an
underspecified question may retrieve evidence for the wrong
document, article, agency, or time period. Therefore, our
system treats \texttt{clarify} as a valid routing action
rather than an exceptional failure case.

\subsection{Dataset and Benchmark Construction}

Recent benchmark papers emphasize that a dataset contribution
should report not only final model scores, but also source
provenance, generation pipeline, label definition, quality
analysis, evaluation protocol, and scope limitations. MIST,
although developed for a different modality and task, is an
example of this reporting style because it describes the data
generation pipeline, statistics, quality analysis, and
task-specific metrics \cite{vu2026mist}. We follow this style
lightly for Vietnamese legal QA: the present dataset is
organized around data sources, construction rules, routing
labels, split protocol, and limitations, but it is not compared
as a method against MIST.

\subsection{Vietnamese NLP and Legal QA}

Vietnamese NLP introduces additional challenges because word
boundaries are not always marked by whitespace and many
entity names consist of multi-syllable expressions.
VnCoreNLP provides Vietnamese word segmentation,
part-of-speech tagging, named entity recognition, and
dependency parsing \cite{vu2018vncorenlp}, while PhoBERT
provides pretrained Vietnamese representations for
downstream NLP tasks \cite{nguyen2020phobert}. Legal NLP
further requires domain-specific reasoning over formal
documents, citations, and regulatory relations
\cite{chalkidis2020legalbert,zhong2020nlp_law}. The
proposed work differs from general-purpose RAG systems by
placing reasoning-aware routing at the center of Vietnamese
legal question answering.

\subsection{Comparative Discussion}
\label{subsec:related_comparison}

Table~\ref{tab:related_comparison_short} summarizes the
relationship between prior research lines and the proposed
system. Existing RAG methods provide a strong foundation for
grounding LLM outputs in retrieved passages, but they mainly
operate over unstructured text chunks and do not explicitly
encode legal relations. Graph-based retrieval methods
improve relational evidence access, but they often assume
that graph retrieval is the main retrieval mechanism instead
of deciding when graph traversal is necessary. Adaptive
routing methods demonstrate that inference paths can be
selected dynamically, but most of them focus on routing
among model sizes, retrieval depths, or QA strategies rather
than routing among dense retrieval, graph traversal, hybrid
reasoning, and clarification. Ambiguity-aware QA methods
show the importance of rewriting or clarifying
underspecified questions, but clarification is rarely
integrated as a retrieval route in Hybrid GraphRAG.

The proposed system differs from these lines by treating
route selection as a reasoning-aware decision before
retrieval. In Vietnamese legal QA, this design is necessary
because the same surface form may correspond to different
legal reasoning requirements. A query asking about a named
article may be answered by dense retrieval, whereas a query
asking whether a decree amends another document requires
graph traversal. A query that combines article content with
amendment or effective-date relations may require hybrid
reasoning. If the legal target is missing, the safest action
is clarification rather than retrieval.

\begin{table}[!htbp]
\centering
\caption{Summary of related research lines.}
\label{tab:related_comparison_short}
\scriptsize
\setlength{\tabcolsep}{2.5pt}
\renewcommand{\arraystretch}{1.08}
\begin{tabularx}{\columnwidth}{L{1.45cm}Y L{1.35cm}}
\toprule
\textbf{Line} &
\textbf{Limitation for this work} &
\textbf{Role} \\
\midrule

Dense RAG
\cite{lewis2020rag,karpukhin2020dpr,guu2020realm} &
Retrieves relevant text but does not explicitly model legal
relations such as amendment, repeal, reference, or effective
date. &
Dense route \\

Adv. RAG
\cite{asai2023selfrag,jiang2023flare,sarthi2024raptor} &
Improves retrieval control or hierarchy, but does not route
among dense, graph, hybrid, and clarify actions. &
Control \\

Graph RAG
\cite{edge2024graphrag,wang2024kgp,li2024graphreader,
gutierrez2024hipporag} &
Supports relational retrieval, but graph traversal can be
unnecessary for simple legal lookup questions. &
Graph route \\

Routing
\cite{jeong2024adaptive,chen2023frugalgpt,ong2024routellm} &
Routes models or retrieval depths, but not legal evidence
mechanisms and clarification decisions. &
Router \\

Ambiguity QA
\cite{min2020ambigqa,elgohary2019canard} &
Handles ambiguous questions, but clarification is rarely
integrated as a Hybrid GraphRAG route. &
Clarify \\

VN Legal NLP
\cite{nguyen2020phobert,vu2018vncorenlp,chalkidis2020legalbert,
zhong2020nlp_law} &
Provides language or legal NLP foundations, but not
reasoning-aware retrieval routing. &
Features \\

\textbf{Ours} &
Selects the evidence mechanism according to reasoning,
ambiguity, and cost constraints. &
4-way route \\

\bottomrule
\end{tabularx}
\end{table}

\FloatBarrier

% ============================================================
\section{Methodology}
\label{sec:methodology}
% ============================================================

\subsection{Problem Definition}
\label{subsec:problem_definition}

Let $\mathcal{C}=\{D_1,D_2,\ldots,D_N\}$ denote the
Vietnamese legal corpus, where $D_i$ is the $i$-th legal
document and $N$ is the number of documents. Each document
contains legal text and metadata such as title, legal type,
issuing authority, source identifier, article structure, and
available temporal information. From $\mathcal{C}$, the
system builds two complementary knowledge sources. The first
source is a dense vector index
$\mathcal{I}_v=\{(c_j,\mathbf{e}_j)\}_{j=1}^{M}$, where
$c_j$ is the $j$-th legal text chunk, $\mathbf{e}_j\in
\mathbb{R}^{d}$ is its $d$-dimensional embedding, and $M$ is
the number of indexed chunks. The second source is a legal
knowledge graph $\mathcal{G}=(\mathcal{V},\mathcal{E})$,
where $\mathcal{V}$ is the set of graph nodes and
$\mathcal{E}$ is the set of directed typed edges. A node
$v\in\mathcal{V}$ may represent a legal document, legal
article, vector chunk, agency, or legal concept. An edge
$(u,r,v)\in\mathcal{E}$ represents a legal relation $r$ from
source node $u$ to target node $v$, for example
\textsc{HAS\_ARTICLE}, \textsc{BELONGS\_TO},
\textsc{REFERENCES}, \textsc{AMENDS}, \textsc{REPEALS}, or
\textsc{REGULATES\_CONCEPT}.

At inference time, the input is the current user query $q_t$
and an optional conversation history
$H=\{(q_1,a_1),\ldots,(q_{t-1},a_{t-1})\}$. Here, $t$
denotes the current dialogue turn, $q_i$ denotes the user
query at turn $i$, and $a_i$ denotes the corresponding
system answer. If the current query is the first turn, then
$H=\emptyset$. The router must choose one action from the
route set
\begin{equation}
\mathcal{D}=\{d_v,d_g,d_h,d_c\},
\label{eq:route_set}
\end{equation}
where $d_v$ denotes \texttt{dense\_retrieval}, $d_g$ denotes
\texttt{graph\_traversal}, $d_h$ denotes
\texttt{hybrid\_reasoning}, and $d_c$ denotes
\texttt{clarify}. The routing function is defined as
\begin{equation}
\rho:(q_t,H)\rightarrow d^*, \qquad d^*\in\mathcal{D},
\label{eq:routing_function}
\end{equation}
where $d^*$ is the final selected route.

If $d^*\neq d_c$, the system retrieves an evidence set
$E_{d^*}$ and generates a grounded answer. For the three
retrieval routes, the evidence set is
\begin{equation}
E_{d^*}=\begin{cases}
E_v=\mathrm{TopK}(q_t,\mathcal{I}_v,k), & d^*=d_v,\\
E_g=\mathrm{Traverse}(q_t,\mathcal{G},h), & d^*=d_g,\\
E_v\cup E_g, & d^*=d_h,
\end{cases}
\label{eq:evidence_set}
\end{equation}
where $k$ is the dense retrieval depth and $h$ is the maximum
graph traversal depth. The generator then produces
\begin{equation}
a_t=\mathrm{Generate}(q_t,H,E_{d^*}).
\label{eq:answer_generation}
\end{equation}
If $d^*=d_c$, the system does not retrieve evidence and
instead returns
\begin{equation}
c_t=\mathrm{Clarify}(q_t,H),
\label{eq:clarification_generation}
\end{equation}
where $c_t$ asks the user to specify the missing legal
document, article, entity, time period, or conversational
referent. This distinction is important in legal QA because
an answer grounded in the wrong legal target can be more
harmful than asking a follow-up question.

\subsection{Dataset Construction}
\label{subsec:dataset_description}

Following the benchmark-reporting style discussed in
Section~\ref{sec:related_work}, we describe the dataset by
its provenance, construction pipeline, label rules, split
protocol, and limitations rather than only by final counts.
The routing dataset is self-constructed from processed
Vietnamese legal sources. The first source contains legal
documents stored as JSONL records with document identifiers,
titles, document types, issuing authorities, and normalized
legal text. The second source is the Vietnamese electronic
codification system, originally organized by theme, topic,
and article. Because retrieval and graph construction require
stable evidence units, codification records are flattened
into article-level objects with fields such as
\texttt{doc\_id}, \texttt{title}, \texttt{theme},
\texttt{topic}, \texttt{content\_markdown}, and source
metadata.

The construction pipeline has four stages. First, legal
documents are normalized and segmented into article-level
records. Second, QA samples are linked to one or more legal
articles through \texttt{article\_key}. Third, these article
keys are parsed into a structured \texttt{relevant\_articles}
list, joined with the article corpus to create
\texttt{gold\_context}, and converted into
\texttt{supporting\_facts}. Finally, route labels are derived
from the supporting evidence structure and a strict
train/dev/test split is produced.

Each QA sample contains a Vietnamese legal question, a
reference answer, an evidence text, a legal source identifier,
and an \texttt{article\_key} connecting the sample to one or
more supporting legal articles. The field
\texttt{article\_key} follows the format
\texttt{DOC\_NUMBER::ARTICLE\_ID}; for example,
\texttt{49/VBHN-VPQH::Điều 75}. Low-quality samples are
removed with deterministic filters: unsupported route labels
are excluded, questions shorter than $12$ characters or
answers shorter than two characters are discarded, duplicate
questions are removed, and placeholder answers such as
\textit{không có trong ngữ cảnh}, \textit{không có thông tin},
or \textit{đang cập nhật nội dung} are filtered out.

Routing labels are derived from evidence structure. A sample
is labeled \texttt{dense\_retrieval} when it can usually be
answered from a single article or direct textual evidence. It
is labeled \texttt{graph\_traversal} when it requires
multiple articles from the same legal document, indicating
intra-document structural or relational reasoning. It is
labeled \texttt{hybrid\_reasoning} when the supporting
evidence spans multiple documents, because these cases often
require both semantic retrieval and relation-aware graph
reasoning. These labels are useful for system development,
but they are metadata-derived rather than independently
human-annotated; this limitation is discussed in
Section~\ref{sec:limitations}.

The strict routing dataset contains $1{,}126$ samples with
three observed retrieval labels. It is split into train,
development, and test sets using stratified sampling with a
fixed random seed. To reduce leakage, the test set is sampled
first and is never used for router training or threshold
tuning; the development set is then sampled from the
remaining pool, and the remaining samples form the training
set.
Table~\ref{tab:dataset} reports the split statistics.

\begin{table}[!htbp]
\centering
\caption{Vietnamese legal QA routing dataset split by route label.}
\label{tab:dataset}
\small
\begin{tabular}{lrrrr}
\toprule
\textbf{Split} & \textbf{Dense} & \textbf{Graph} & \textbf{Hybrid} & \textbf{Total} \\
\midrule
Train & 356 & 54 & 36 & 446 \\
Dev & 50 & 20 & 10 & 80 \\
Test & 300 & 150 & 150 & 600 \\
\midrule
\textbf{Total} & 706 & 224 & 196 & 1,126 \\
\bottomrule
\end{tabular}
\end{table}

\begin{figure}[!htbp]
\centering
\begin{tikzpicture}
\begin{axis}[
    ybar,
    width=0.95\columnwidth,
    height=5.0cm,
    ymin=0,
    ylabel={Number of samples},
    symbolic x coords={Dense, Graph, Hybrid},
    xtick=data,
    bar width=16pt,
    nodes near coords,
    nodes near coords style={font=\scriptsize},
    tick label style={font=\scriptsize},
    label style={font=\scriptsize},
    enlarge x limits=0.25
]
\addplot coordinates {(Dense,706) (Graph,224) (Hybrid,196)};
\end{axis}
\end{tikzpicture}
\caption{Route-label distribution of the strict Vietnamese legal QA routing dataset. The distribution is skewed toward dense retrieval, so Macro-F1 and per-class analysis are needed in addition to accuracy.}
\label{fig:dataset_distribution}
\end{figure}

Because the strict split does not contain standalone
\texttt{clarify} training labels, clarification is evaluated
on a separate constructed ambiguity benchmark. This synthetic
stress-test set contains $234$ queries: $156$ ambiguous
queries labeled \texttt{clarify} and $78$ answerable queries
labeled as one of the three retrieval routes. The ambiguous
cases are organized into four types: incomplete context
($39$), missing entity ($39$), unresolved pronoun reference
($42$), and multiple interpretation ($36$). This separation is
deliberate: the strict split evaluates retrieval-mechanism
selection, whereas the ambiguity benchmark evaluates whether
the router can avoid unsafe answering when the legal target is
underspecified. Since the ambiguity set is template-generated
rather than independently human-annotated, it should be
interpreted as a controlled stress test, not as an external
public benchmark.

After Phase~3, we also construct a $160$-query
conversation-aware ambiguity stress test. This diagnostic set
contains eight equally sized categories:
\texttt{answerable\_with\_history},
\texttt{clarify\_without\_history},
\texttt{irrelevant\_history},
\texttt{conflicting\_history}, \texttt{missing\_entity},
\texttt{multi\_interpretation}, \texttt{clear\_dense\_control},
and \texttt{clear\_graph\_or\_hybrid\_control}. Its purpose is
not to replace the strict $600$-query end-to-end benchmark,
but to test whether the router uses conversation history
correctly: a resolved referent should allow retrieval, while
absent, irrelevant, or conflicting history should preserve or
force clarification. The benchmark is template-controlled and
should therefore be interpreted as a stress test for router
behavior rather than as a production distribution.
Table~\ref{tab:dataset_protocol} summarizes the construction
protocol for the strict legal QA set and the original
ambiguity stress test.

\begin{table}[!htbp]
\centering
\caption{Dataset construction protocol for the primary
evaluation sets used in this study.}
\label{tab:dataset_protocol}
\scriptsize
\setlength{\tabcolsep}{2.0pt}
\renewcommand{\arraystretch}{1.08}
\begin{tabularx}{\columnwidth}{L{1.8cm}YY}
\toprule
\textbf{Item} & \textbf{Strict legal QA set} &
\textbf{Ambiguity stress test} \\
\midrule
Purpose &
Evaluate retrieval-route selection and end-to-end QA quality. &
Evaluate whether the router asks for clarification before
unsafe retrieval. \\
Source &
Processed Vietnamese legal QA samples linked to legal
articles by \texttt{article\_key}. &
Deterministic templates for underspecified legal questions
plus clear answerable controls. \\
Labels &
Dense, graph, and hybrid labels derived from evidence
structure. &
\texttt{clarify} positives and dense/graph/hybrid negative
controls. \\
Size &
$1{,}126$ total samples; $600$ held out for strict testing. &
$234$ total queries; $156$ clarify and $78$ answerable controls. \\
Quality control &
Deduplication, short-answer filtering, placeholder-answer
filtering, and held-out test-first split. &
Fixed generation rules, balanced control routes, and explicit
ambiguity-type labels. \\
Scope &
Tests retrieval strategy selection, not standalone
clarification. &
Controlled stress test; requires human review before being
treated as a production benchmark. \\
\bottomrule
\end{tabularx}
\end{table}

\subsection{System Pipeline}
\label{subsec:system_pipeline}

Figure~\ref{fig:system_pipeline} summarizes the proposed
pipeline. The input $(q_t,H)$ is first converted into
query-level reasoning and ambiguity features. In the
conversation-aware implementation, a deterministic
\texttt{HistoryResolver} first produces $r_H$ so that the
feature extractor can distinguish resolved, missing,
irrelevant, and conflicting history. Stage~1
predicts an initial route using a fast feature-based model.
Stage~2 is invoked only when the trigger function detects low
confidence, ambiguity, or high reasoning demand. The final
route then controls whether the system executes dense
retrieval, graph traversal, hybrid retrieval, or
clarification.

\begin{figure*}[p]
\centering
\resizebox{0.94\textwidth}{!}{
\begin{tikzpicture}[
    font=\small,
    >=Latex,
    node distance=13mm and 15mm,
    main/.style={rectangle,rounded corners=2pt,draw,align=center,minimum height=12mm,minimum width=32mm,fill=gray!8},
    data/.style={rectangle,draw,align=center,minimum height=12mm,minimum width=32mm,fill=blue!5},
    route/.style={rectangle,rounded corners=2pt,draw,align=center,minimum height=15mm,minimum width=39mm,fill=green!8},
    output/.style={rectangle,rounded corners=2pt,draw,align=center,minimum height=12mm,minimum width=35mm,fill=purple!8},
    example/.style={rectangle,rounded corners=3pt,draw,align=center,minimum height=34mm,minimum width=52mm,fill=orange!10},
    note/.style={rectangle,rounded corners=3pt,draw,align=left,text width=0.88\textwidth,fill=yellow!8,inner sep=8pt},
    arrow/.style={->, thick},
    exarrow/.style={->, very thick, dashed}
]
\node[data] (input) {Input\\query $q_t$\\history $H$};
\node[main, right=of input] (features) {Feature extraction\\$\mathbf{f}(q_t,H)$\\$\phi_{\mathrm{amb}}(q_t,H)$, $\xi(q_t)$};
\node[main, right=of features] (stage1) {Stage~1 router\\XGBoost\\$(d_1,s_1)$\\{\footnotesize $d_1\in\{d_v,d_g,d_h\}$}};
\node[main, right=of stage1] (trigger) {Trigger check\\$\Gamma(q_t,H,r_H)$};
\node[main, above=16mm of trigger] (stage2) {Stage~2 verifier\\LLM\\$f_2(q_t,H,d_1)$\\{\footnotesize $d^*\in\mathcal{D}$, incl. $d_c$}};
\node[main, right=of trigger] (finalroute) {Final route\\$d^* \in \mathcal{D}$};
\node[route, right=of finalroute] (execution) {Route execution\\[1mm]$d_v$: dense retrieval\\$d_g$: graph traversal\\$d_h$: hybrid reasoning\\$d_c$: clarify};
\node[main, right=of execution] (evidence) {Evidence assembly\\or clarification\\$E_{d^*}$ or $c$};
\node[main, below=15mm of evidence, minimum width=42mm, fill=red!5] (fallback) {Evidence guard\\fallback if\\uninformative\\dense$\leftrightarrow$graph\\hybrid$\rightarrow$vector};
\node[output, right=of evidence] (out) {Output\\answer $a$\\or question $c$};
\node[data, below=16mm of execution, xshift=-28mm] (dense) {Dense index\\$\mathcal{I}_v$};
\node[data, below=16mm of execution, xshift=28mm] (kg) {Legal KG\\$\mathcal{G}$};
\draw[arrow] (input) -- (features);
\draw[arrow] (features) -- (stage1);
\draw[arrow] (stage1) -- (trigger);
\draw[arrow] (trigger) -- node[above]{no} (finalroute);
\draw[arrow] (trigger) -- node[left]{yes} (stage2);
\draw[arrow] (stage2.east) -| (finalroute.north);
\draw[arrow] (finalroute) -- (execution);
\draw[arrow] (execution) -- (evidence);
\draw[arrow] (evidence) -- (out);
\draw[arrow, dashed] (evidence.south) -- (fallback.north);
\draw[arrow, dashed] (fallback.west) -- ++(-5mm,0) |- (execution.south east);
\draw[arrow] (dense.north) -- ++(0,7mm) -| (execution.south);
\draw[arrow] (kg.north) -- ++(0,7mm) -| (execution.south);
\node[main, below=24mm of features, minimum width=86mm] (routeset) {Route set $\mathcal{D}=\{d_v,d_g,d_h,d_c\}$\\$d_v$: dense retrieval,\quad $d_g$: graph traversal\\$d_h$: hybrid reasoning,\quad $d_c$: clarification};
\draw[arrow, dashed] (routeset.east) -| (finalroute.south);
\node[draw,dashed,rounded corners,fit=(features)(stage1)(trigger)(stage2)(finalroute),inner sep=5mm,label={[font=\small]above:Reasoning-aware adaptive router}] {};
\node[draw,dashed,rounded corners,fit=(execution)(dense)(kg)(evidence)(fallback),inner sep=5mm,label={[font=\small]below:Retrieval and response execution}] {};
\node[example, below=58mm of input, minimum width=54mm] (exdense) {\textbf{1. Dense route $d_v$}\\[2mm]$q$: ``\textit{Điều kiện cấp}\\\textit{giấy phép xây dựng là gì?}''\\[2mm]Signal:\\direct legal lookup\\[1mm]Action:\\retrieve top-$k$ chunks\\from $\mathcal{I}_v$};
\node[example, right=of exdense, minimum width=54mm] (exgraph) {\textbf{2. Graph route $d_g$}\\[2mm]$q$: ``\textit{Văn bản nào sửa đổi}\\\textit{Nghị định 100/2019?}''\\[2mm]Signal:\\amendment relation\\[1mm]Action:\\traverse \textsc{AMENDS}\\or \textsc{REFERENCES} edges};
\node[example, right=of exgraph, minimum width=54mm] (exhybrid) {\textbf{3. Hybrid route $d_h$}\\[2mm]$q$: ``\textit{Mức phạt hiện nay}\\\textit{theo văn bản sửa đổi}\\\textit{là bao nhiêu?}''\\[2mm]Signal:\\relation + article text\\[1mm]Action:\\combine $E_v \cup E_g$};
\node[example, right=of exhybrid, minimum width=54mm] (exclarify) {\textbf{4. Clarify route $d_c$}\\[2mm]$q$: ``\textit{Văn bản đó}\\\textit{còn hiệu lực không?}''\\$H=\emptyset$\\[2mm]Signal:\\unresolved referent\\[1mm]Action:\\ask clarification};
\node[draw,dashed,rounded corners,fit=(exdense)(exgraph)(exhybrid)(exclarify),inner sep=7mm,label={[font=\small]above:Concrete Vietnamese examples for four routing outcomes}] {};
\node[note, below=18mm of exgraph, xshift=39mm] (exnote) {\textbf{Interpretation.} Simple lookup questions can be handled by dense retrieval; relation-oriented questions benefit from graph traversal; relation-plus-content questions require hybrid evidence; and underspecified questions should be clarified before retrieval.};
\end{tikzpicture}}
\caption{Overview of the proposed Reasoning-aware Adaptive Router for Vietnamese Legal Hybrid GraphRAG. The figure also explains the notation used in the paper: $q_t$ is the current query, $H$ is the conversation history, $\mathbf{f}$ is the extracted feature vector, $\phi_{\mathrm{amb}}$ is the ambiguity score, $\xi$ is the multi-hop score, $(d_1,s_1)$ is the Stage~1 route and confidence, and $d^*$ is the final route. The dashed fallback guard indicates the implementation-level retry policy, detailed in Table~\ref{tab:route_fallback_policy}.}
\label{fig:system_pipeline}
\end{figure*}

\subsection{Running Example}
\label{subsec:running_example}

This subsection illustrates how the formal routing definition
and Figure~\ref{fig:system_pipeline} are used for different
Vietnamese legal questions. Consider the query
$q_1=$ ``\textit{Điều kiện để được cấp giấy phép xây dựng là
gì?}''. This query expresses a direct legal information need.
It asks for conditions that are usually stated in one or a
small number of legal provisions, and it does not explicitly
mention amendment, repeal, replacement, cross-document
reference, or effective-date reasoning. The feature extractor
therefore assigns low ambiguity and low multi-hop signals. In
this case, Stage~1 is expected to predict the dense route
$d_v$ with high confidence, so Stage~2 is bypassed and the
system retrieves relevant chunks from the dense index
$\mathcal{I}_v$. The evidence set is
$E_{d^*}=E_v=\mathrm{TopK}(q_1,\mathcal{I}_v,k)$, and the
answer generator produces a grounded answer from the retrieved
legal text.

In contrast, consider $q_2=$ ``\textit{Nghị định nào sửa đổi
quy định về xử phạt vi phạm hành chính trong lĩnh vực đất
đai, và điều khoản hiệu lực của nghị định đó được quy định ở
đâu?}''. This query contains explicit relation signals such
as \textit{sửa đổi} and \textit{hiệu lực}. It also requires
connecting at least two legal objects: the document being
amended and the decree that introduces the amendment. Moreover,
the phrase asking where the validity clause is specified
requires locating an article or provision inside the modifying
decree. The feature extractor therefore assigns stronger
legal-relation and multi-hop signals. If the answer can be
obtained mainly from legal relations, the router selects
$d_g$ and retrieves graph paths from $\mathcal{G}$. If both
relation paths and article text are needed, the router selects
$d_h$ and constructs $E_{d^*}=E_v\cup E_g$. This example shows
why graph traversal should be available but should not be
applied to every legal query.

Finally, consider $q_3=$ ``\textit{Văn bản đó còn hiệu lực
không?}''. The phrase \textit{văn bản đó} is an unresolved
referent when the conversation history $H$ does not identify a
unique legal document. If the system retrieves evidence
immediately, it may answer about the wrong document or time
period. Therefore, the ambiguity score
$\phi_{\mathrm{amb}}(q_3,H)$ should be high, the Stage~2
verifier may be triggered, and the final route should be
$d_c$. Instead of producing an answer, the system returns a
clarification question such as ``\textit{Bạn muốn hỏi hiệu lực
của văn bản nào?}''. This behavior is important in legal QA
because asking for clarification is safer than generating an
answer from an incorrectly inferred legal target.

\subsection{Legal Knowledge Base Construction}
\label{subsec:legal_kb_construction}

The retrieval backend uses \path{microsoft/Harrier-OSS-v1-0.6B}
embeddings and represents $199{,}530$ graph-linked vector
chunks as Neo4j \textsc{VectorChunk} nodes. Dense retrieval is
appropriate for direct questions whose answers are stated in
legal text, such as definitions, conditions, penalties, or
procedures. The graph backend stores explicit relations among
legal documents, articles, chunks, and legal concepts. This
graph is necessary because many legal questions require
following explicit legal links, such as article membership,
references, amendments, repeal, and concept regulation.

Table~\ref{tab:graph_stats} reports the current graph after
the Pháp điển migration and vector-chunk linking. After the
Pháp điển migration, the most important schema change is that
Pháp điển article nodes are now exposed as
\textsc{LegalArticle} nodes with content and connected to
parent legal documents through \textsc{HAS\_ARTICLE}. This
change is important because article-level retrieval is the
natural evidence unit for legal QA.

\begin{table}[!htbp]
\centering
\caption{Statistics of the enriched Vietnamese legal graph $\mathcal{G}$ after graph migration. Vector chunks are reported using the post-migration graph-linked \textsc{VectorChunk} count.}
\label{tab:graph_stats}
\small
\begin{tabular}{lr}
\toprule
\textbf{Graph item} & \textbf{Count} \\
\midrule
Nodes $|\mathcal{V}|$ & 419,251 \\
Edges $|\mathcal{E}|$ & 1,239,542 \\
Legal documents & 219,439 \\
Legal articles & 70,347 \\
Neo4j \textsc{VectorChunk} nodes & 199,530 \\
Legal concepts & 10 \\
Average degree & 5.91 \\
Isolated nodes & 3,102 \\
\bottomrule
\end{tabular}
\end{table}

\begin{table}[!htbp]
\centering
\caption{Main relation types in the enriched legal graph.}
\label{tab:graph_relations}
\small
\begin{tabular}{lr}
\toprule
\textbf{Relation type} & \textbf{Count} \\
\midrule
\textsc{REFERENCES} & 769,000 \\
\textsc{HAS\_VECTOR\_CHUNK} & 199,802 \\
\textsc{BELONGS\_TO} & 199,530 \\
\textsc{HAS\_ARTICLE} & 70,347 \\
\textsc{REGULATES\_CONCEPT} & 249 \\
\textsc{NEXT\_ARTICLE} & 239 \\
\textsc{CROSS\_REFERENCES} & 133 \\
\textsc{MENTIONS\_DOC} & 84 \\
\textsc{INTRA\_DOC\_REFERENCE} & 83 \\
\textsc{AMENDS} & 27 \\
\textsc{GUIDES} & 25 \\
\textsc{REPEALS} & 19 \\
\textsc{IMPLEMENTS} & 3 \\
\bottomrule
\end{tabular}
\end{table}

These statistics show that the graph is large and well linked
at the document and chunk level, while curated legal-effect
relations such as \textsc{AMENDS}, \textsc{REPEALS}, and
\textsc{GUIDES} remain sparse. The \textsc{HAS\_ARTICLE}
count of $70{,}347$ reflects the graph migration step that
promoted Pháp điển article nodes from \textsc{LegalDoc} to
\textsc{LegalArticle}; the \textsc{BELONGS\_TO} relation
connects vector chunks back to their source legal documents,
which is important for hybrid retrieval and evidence
serialization.

The graph backend is therefore treated as a selective
reasoning resource rather than as the default retrieval
mechanism. Relational structure can improve multi-hop legal
evidence access, but uncontrolled graph expansion may also
introduce broad or noisy context. In the current graph,
generic \textsc{REFERENCES} edges are abundant, whereas
curated legal-effect relations such as \textsc{AMENDS},
\textsc{REPEALS}, and \textsc{GUIDES} remain sparse. As a
result, graph traversal is valuable for relation-heavy
queries but can be unnecessarily costly or diffuse for simple
article-level lookup questions. This observation directly
motivates adaptive routing instead of always applying
GraphRAG.

\subsection{Reasoning-aware Feature Extraction}
\label{subsec:feature_extraction}

For each input $(q,H)$, the router extracts a feature vector
$\mathbf{f}(q,H)$ and two auxiliary scores: the ambiguity
score $\phi_{\mathrm{amb}}(q,H)$ and the multi-hop score
$\xi(q)$. The feature vector combines statistical features,
Vietnamese legal-reference patterns, entity features,
question-type cues, relation cues, and context-dependence
signals. Statistical features include query length, token
count, named-entity count, entity density, and the number of
legal identifiers. Question-type features encode Vietnamese
question forms such as \textit{gì}, \textit{khi nào},
\textit{ở đâu}, \textit{bao nhiêu}, and \textit{như thế nào}.
Legal-relation features detect phrases such as \textit{sửa
đổi}, \textit{bãi bỏ}, \textit{thay thế}, \textit{có hiệu
lực}, and \textit{được quy định tại}. Multi-hop features
capture comparison markers, conditional clauses, and signals
that more than one article or document may be required.
Ambiguity features detect unresolved pronouns, vague legal
targets, missing entities, missing temporal scopes, and
queries that depend on absent conversation history.

Phase~3 adds a deterministic \texttt{HistoryResolver} before
ambiguity scoring. It extracts legal referent candidates from
the conversation history and returns
\begin{equation}
r_H=\mathrm{Resolve}(q_t,H),
\label{eq:history_resolver}
\end{equation}
where $r_H$ contains \texttt{candidate\_referents}, an optional
\texttt{resolved\_referent}, a
\texttt{history\_resolution\_status}, and a
\texttt{history\_resolution\_confidence}. The possible status
values are \texttt{resolved}, \texttt{no\_history},
\texttt{irrelevant\_history}, \texttt{conflicting\_history},
and \texttt{not\_needed}. This signal prevents the router from
treating every non-empty history as useful context. If a
query such as ``\textit{Văn bản đó còn hiệu lực không?}''
has a unique document referent in history, the resolved
referent can unblock graph traversal. If history is absent,
irrelevant, or contains multiple competing referents, the same
surface form remains ambiguous and can be routed to
clarification.

\subsection{Two-stage Adaptive Router}
\label{subsec:two_stage_router}

Stage~1 is an XGBoost classifier $f_1$ trained on the strict
three-class routing data. It maps the extracted feature
vector to an initial route and confidence score:
\begin{equation}
(d_1,s_1)=f_1(\mathbf{f}(q,H)),
\label{eq:stage1_router}
\end{equation}
where $d_1$ is the initial route and $s_1\in[0,1]$ is the
confidence score. Because the Stage~1 training split contains
only \texttt{dense\_retrieval}, \texttt{graph\_traversal},
and \texttt{hybrid\_reasoning}, Stage~1 is not expected to
learn the \texttt{clarify} class. This is an architectural
property of the current training setup rather than an
implementation bug.

Stage~2 is an LLM-based verifier $f_2$ invoked only for
selected cases. Let $b=\phi_{\mathrm{amb}}(q,H)$ be the
ambiguity score. The trigger indicators are
\begin{align}
\gamma_{\mathrm{conf}} &= \mathbb{I}[s_1 < \tau_{\mathrm{conf}}], \\
\gamma_{\mathrm{amb}} &= \mathbb{I}[b > \tau_{\mathrm{force}}], \\
\gamma_{\mathrm{hop}} &= \mathbb{I}[d_1 \neq d_g \wedge \xi(q)=2].
\end{align}
The trigger function is
\begin{equation}
\Gamma(q,H,r_H)=\mathbb{I}[\gamma_{\mathrm{conf}}\vee
\gamma_{\mathrm{amb}}\vee\gamma_{\mathrm{hop}}].
\label{eq:trigger_function}
\end{equation}
The final route is then
\begin{equation}
d^*=\begin{cases}
f_2(q,H,d_1), & \Gamma(q,H,r_H)=1,\\
d_1, & \Gamma(q,H,r_H)=0.
\end{cases}
\label{eq:stage2_router}
\end{equation}
The current default implementation uses
$\tau_{\mathrm{conf}}=0.50$ for low-confidence triggering,
$\tau_{\mathrm{force}}=0.60$ for ambiguity-based Stage~2
forcing, and a separate clarification override threshold
$\tau_{\mathrm{clar}}=0.80$. The ablation in
Section~\ref{subsec:ablation_results} shows that
$\tau_{\mathrm{conf}}$ mainly changes how often Stage~2 is
considered, while $\tau_{\mathrm{clar}}$ has a stronger effect
on false clarification in non-ambiguous routing data.
The Phase~3 implementation additionally uses $r_H$: resolved
history can block unnecessary clarification and restore a
retrieval route, while \texttt{no\_history},
\texttt{irrelevant\_history}, or
\texttt{conflicting\_history} can force Stage~2 or
clarification when the query contains a contextual reference.
Table~\ref{tab:stage2_policy} summarizes the implemented
Stage~2 verification policy.

\begin{table}[!htbp]
\centering
\caption{Stage~2 verification policy used by the implemented two-stage router.}
\label{tab:stage2_policy}
\scriptsize
\setlength{\tabcolsep}{1.8pt}
\renewcommand{\arraystretch}{0.96}
\begin{tabularx}{\columnwidth}{L{2.15cm}L{3.35cm}Y}
\toprule
\textbf{Signal} & \textbf{Condition} & \textbf{Purpose} \\
\midrule
Low confidence &
$s_1 < \tau_{\mathrm{conf}}$ &
Check uncertain Stage~1 decisions. \\
Ambiguity risk &
$b \ge \tau_{\mathrm{force}}$ &
Avoid retrieval when the legal target is underspecified. \\
Resolved history &
$r_H.\mathrm{status}=\texttt{resolved}$ &
Use the referent and avoid unnecessary clarification. \\
Unresolved history &
$r_H.\mathrm{status}$ is no, irrelevant, or conflicting &
Force verification or clarification. \\
Dense fast path &
High-confidence dense, low ambiguity, no strong legal signal &
Skip Stage~2 for clear lookup queries. \\
Dense legal signal &
Dense route with legal-reference, graph, authority, or effect signals &
Catch cases needing graph or hybrid retrieval. \\
Reasoning signal &
Multi-hop, cross-document, complex, or relation-heavy features &
Verify relation-heavy legal questions. \\
Clarify sanity check &
Predicted \texttt{clarify} with low ambiguity score &
Reduce false clarification on answerable questions. \\
\bottomrule
\end{tabularx}
\end{table}

\begin{algorithm}[!htbp]
\caption{Two-stage Adaptive Router}
\label{alg:two_stage_router}
\begin{algorithmic}[1]
\Require Query $q$, history $H$, thresholds $\tau_{\mathrm{conf}},\tau_{\mathrm{force}},\tau_{\mathrm{clar}}$
\Ensure Final route $d^*\in\mathcal{D}$
\State $r_H\leftarrow\mathrm{Resolve}(q,H)$
\State Compute $\mathbf{f}(q,H,r_H)$, $\phi_{\mathrm{amb}}(q,H,r_H)$, and $\xi(q)$
\State $(d_1,s_1)\leftarrow f_1(\mathbf{f}(q,H,r_H))$
\State Compute $\Gamma(q,H,r_H)$ using Eq.~\eqref{eq:trigger_function}
\If{$\Gamma(q,H,r_H)=1$}
  \State $d^*\leftarrow f_2(q,H,d_1)$
\Else
  \State $d^*\leftarrow d_1$
\EndIf
\State \Return $d^*$
\end{algorithmic}
\end{algorithm}

Algorithm~\ref{alg:two_stage_router} presents the routing
logic in simplified form. The deployed implementation adds
several practical heuristics beyond the three-indicator
trigger: a high-confidence dense fast path that skips
Stage~2 when the query has no legal-reference or
relational signal, a reasoning-signal check that can
escalate queries with strong legal features even at
moderate Stage~1 confidence, and a rescue-only override
policy for Stage~2 that permits route changes only under
narrowly defined conditions. These heuristics are
configurable and are designed to reduce unnecessary
Stage~2 invocations on clear queries while preserving
verification for uncertain cases.

\subsection{Route Execution}
\label{subsec:route_execution}

For $d_v$, the system retrieves the top-$k$ chunks from the
dense index $\mathcal{I}_v$ and sends the selected chunks to
the answer generator. For $d_g$, the system links query
entities to graph nodes and traverses relevant relations in
$\mathcal{G}$, returning serialized graph paths and linked
legal evidence. For $d_h$, the system combines semantic text
evidence and graph evidence as $E_v\cup E_g$. For $d_c$, the
system returns a clarification question and does not generate
a final legal answer. This route execution design matches the
main motivation of the project: not every query needs graph
traversal, and not every query should be answered
immediately. Table~\ref{tab:route_fallback_policy} summarizes
the fallback behavior used by the implemented pipeline.

\begin{table}[!htbp]
\centering
\caption{Route execution and fallback policy in the implemented Hybrid GraphRAG pipeline.}
\label{tab:route_fallback_policy}
\scriptsize
\setlength{\tabcolsep}{2.3pt}
\renewcommand{\arraystretch}{1.08}
\begin{tabularx}{\columnwidth}{L{2.2cm}L{3.55cm}Y}
\toprule
\textbf{Route} & \textbf{Primary execution} & \textbf{Fallback / guard} \\
\midrule
\makecell[l]{\texttt{dense}\\\texttt{retrieval}} &
Retrieve from Chroma and generate from dense text chunks. &
If the answer is uninformative, optionally retry graph retrieval. \\
\makecell[l]{\texttt{graph}\\\texttt{traversal}} &
Retrieve graph context through the GraphRAG adapter and serialize legal paths/evidence. &
If graph retrieval or generated evidence is uninformative, fallback to dense retrieval. \\
\makecell[l]{\texttt{hybrid}\\\texttt{reasoning}} &
Merge vector context and graph context, then generate from the combined evidence. &
If hybrid synthesis is uninformative and vector evidence exists, fallback to dense retrieval. \\
\texttt{clarify} &
Generate a follow-up clarification question. &
No retrieval is executed; if LLM clarification fails, use the ambiguity detector's fallback question. \\
\bottomrule
\end{tabularx}
\end{table}

\subsection{Implementation Details}
\label{subsec:implementation_details}

The implementation uses a Chroma vector store for dense
retrieval and Neo4j for graph retrieval. The main Vietnamese
legal vector store is persisted at
\path{data/vector_store/chroma_harrier_oss_0_6b}. Its chunks
are embedded with \path{microsoft/Harrier-OSS-v1-0.6B}. In
the graph, the corresponding chunk layer is represented by
\textsc{VectorChunk} nodes and connected back to source legal
documents through \textsc{BELONGS\_TO} edges. This connection
is useful when hybrid retrieval must serialize both graph
paths and source text evidence.

Stage~1 uses an XGBoost classifier over the project's
hand-crafted reasoning and ambiguity features. The model is
stored as a local router checkpoint and is used for fast
feature-based routing at inference time. Stage~2 is served
through an OpenAI-compatible LLM endpoint; in the reported
experiments, the local configuration uses a Qwen3-compatible
model for structured route verification. The verifier is
prompted to return a route decision, confidence, optional
clarification question, and reasoning trace, which are then
parsed by the two-stage router.

Vietnamese named entity recognition uses
\path{NlpHUST/ner-vietnamese-electra-base} for extracting
candidate legal entities and agencies. Graph retrieval uses
Neo4j full-text indexes over legal documents, legal articles,
and vector chunks. These implementation choices are designed
to keep Stage~1 lightweight while allowing Stage~2 and graph
retrieval to be invoked only when the query requires
additional reasoning or clarification.

\FloatBarrier

% ============================================================
\section{Experimental Evaluation}
\label{sec:eval}
% ============================================================

\subsection{Research Questions}
\label{subsec:research_questions}

The evaluation addresses four research questions. RQ1 asks
whether adaptive routing improves answer quality compared
with always using dense retrieval or always using graph
retrieval. RQ2 asks whether the router selects retrieval
routes accurately on the strict Vietnamese legal QA test set.
RQ3 asks whether Stage~2 improves clarification behavior on
ambiguous queries and what types of ambiguity it can detect.
RQ4 asks whether the latency overhead of Stage~2 justifies
selective verification rather than invoking the LLM verifier
for every query.

\subsection{Experimental Setup}
\label{subsec:experimental_setup}

The strict benchmark contains $600$ held-out legal QA queries
with route labels distributed as $300$ dense, $150$ graph,
and $150$ hybrid. It is used for end-to-end answer-quality
and routing-accuracy evaluation. The ambiguity benchmark
contains $234$ queries and is used only to evaluate
clarification behavior. The two benchmarks are kept separate
because the strict test set is designed to evaluate retrieval
strategy selection, while the ambiguity benchmark is designed
to test whether the system can ask follow-up questions before
unsafe retrieval.

The conversation-aware ambiguity stress test contains $160$
queries and is reported as an additional diagnostic after
Phase~3. It is separated from the strict benchmark because it
evaluates a different behavior: whether conversation history
resolves or fails to resolve contextual references. Therefore,
its metrics are not used to replace the strict end-to-end QA
results in Table~\ref{tab:end_to_end_results}.

\begin{table}[!htbp]
\centering
\caption{End-to-end systems evaluated on the strict Vietnamese legal QA benchmark.}
\label{tab:evaluated_systems}
\scriptsize
\setlength{\tabcolsep}{2.4pt}
\renewcommand{\arraystretch}{1.08}
\begin{tabularx}{\columnwidth}{L{2.65cm}Y}
\toprule
\textbf{System} & \textbf{Description} \\
\midrule
Pure Vector &
Always uses dense retrieval from the Chroma index. This
baseline measures the quality of text-only semantic retrieval. \\
Pure Graph &
Always uses graph traversal in Neo4j. This baseline measures
whether relation-only retrieval is sufficient for legal QA. \\
Single-stage Router &
Uses only the Stage~1 XGBoost route prediction and then
executes the selected dense, graph, hybrid, or clarify action. \\
Two-stage Hybrid &
Uses Stage~1 plus the selective Stage~2 LLM verifier for
uncertain, ambiguous, or reasoning-heavy cases. This is the
proposed full system. \\
\bottomrule
\end{tabularx}
\end{table}

Answer quality is reported with Exact Match (EM) and
token-level F1. Because generated Vietnamese legal answers
are often paraphrastic and longer than the reference answer,
F1 is more informative than EM. Routing quality is reported
with routing accuracy, Stage~2 trigger rate, Stage~2 override
rate, and latency.

\subsection{Metrics}
\label{subsec:metrics}

For route classification, accuracy is the proportion of
correct route decisions. For a route label $\ell$, precision,
recall, and F1 are defined as
\begin{align}
P_\ell &= \frac{TP_\ell}{TP_\ell+FP_\ell},\\
R_\ell &= \frac{TP_\ell}{TP_\ell+FN_\ell},\\
F1_\ell &= \frac{2P_\ell R_\ell}{P_\ell+R_\ell}.
\end{align}
Macro-F1 averages $F1_\ell$ equally over labels, while
Weighted-F1 weights each label by support. Macro-F1 remains
important because the dataset is imbalanced toward dense
retrieval. For clarification, precision measures whether
queries predicted as \texttt{clarify} are truly ambiguous,
recall measures how many ambiguous queries are detected, and
clarify F1 summarizes the trade-off.

\FloatBarrier

% ============================================================
\section{Results and Discussion}
\label{sec:results_discussion}
% ============================================================

\subsection{End-to-end Answer Quality on the Strict Test Set}
\label{subsec:end_to_end_results}

Table~\ref{tab:end_to_end_results} reports the strict
$600$-query benchmark for \textbf{phapdien\_strict} (PD-only).
Pure Vector obtains F1 $=0.4947$,
while Pure Graph obtains F1 $=0.5323$. On this PD corpus,
graph direct lookup substantially improves end-to-end quality for graph
routes; however, the routed router systems still do not exceed pure_graph.
The Single-stage Router attains
F1 $=0.4711$, and the Two-stage Hybrid attains
F1 $=0.4693$ (both below pure_graph, but close to each other on this
template-separable benchmark).
Exact Match (EM) is near zero ($0.0017$) for all systems
because Vietnamese legal answers are typically paraphrastic,
longer than the reference text, and involve normalization
variations. Since EM does not discriminate among systems in
this setting, token-level F1 is used as the primary
answer-quality metric.

\begin{table}[!htbp]
\centering
\caption{End-to-end QA results on the strict Vietnamese legal QA test set. EM is near zero for all systems because legal answers are paraphrastic; F1 is the primary metric.}
\label{tab:end_to_end_results}
\scriptsize
\setlength{\tabcolsep}{2.4pt}
\resizebox{\columnwidth}{!}{%
\begin{tabular}{lrrrr}
\toprule
\textbf{System} & \textbf{EM} & \textbf{F1} & \textbf{Routing Acc.} & \textbf{Latency} \\
\midrule
Pure Vector & 0.0000 & 0.4947 & 0.5000 & 2,139.0 ms \\
Pure Graph & 0.0000 & 0.5323 & 0.2500 & 3,356.7 ms \\
Single-stage Router & 0.0000 & 0.4711 & 1.0000 & 4,530.8 ms \\
Two-stage Hybrid & 0.0000 & 0.4693 & 1.0000 & 4,917.2 ms \\
\bottomrule
\end{tabular}%
}
\end{table}

The routing results demonstrate that the adaptive router is
able to select the correct route on the \textbf{phapdien\_strict}
template benchmark: Single-stage Router and Two-stage Hybrid both
reach routing accuracy $1.0000$. Nevertheless, answer quality on
graph/hybrid routes remains low, which indicates that the limiting
factor is \emph{retrieved evidence quality} rather than routing
correctness. Pure Vector has routing accuracy $0.5000$ because
half of the expected routes are dense; Pure Graph has routing
accuracy $0.2500$ because only one quarter are graph.

To explain the gap despite perfect routing, we observe that
even with 100\% route selection, answer quality is still limited by
retrieval and evidence construction inside each backend.
In particular, graph/hybrid answers remain sensitive to whether the
retrieved neighborhood contains the intended legal article text,
so routed router systems can lag behind the pure\_graph baseline.

Table~\ref{tab:per_class_stage1} reports the offline Stage~1
classifier behavior on the strict test split. This diagnostic
is useful because the strict routing data are imbalanced. The
classifier is strongest on \texttt{dense\_retrieval} and
\texttt{hybrid\_reasoning}, while \texttt{graph\_traversal}
is the hardest class. The offline classifier reaches accuracy
$1.000$ and Macro-F1 $1.000$ on the PD-only strict test
split; this diagnostic is separate
from the deployed end-to-end routing metric in
Table~\ref{tab:end_to_end_results}.

\begin{table}[!htbp]
\centering
\caption{Offline Stage~1 router per-class performance on the strict test split.}
\label{tab:per_class_stage1}
\scriptsize
\setlength{\tabcolsep}{3.0pt}
\begin{tabular}{lrrrr}
\toprule
\textbf{Class} & \textbf{Precision} & \textbf{Recall} & \textbf{F1} & \textbf{Support} \\
\midrule
Dense retrieval & 1.000 & 1.000 & 1.000 & 300 \\
Graph traversal & 1.000 & 1.000 & 1.000 & 150 \\
Hybrid reasoning & 1.000 & 1.000 & 1.000 & 150 \\
\midrule
Macro average & 1.000 & 1.000 & 1.000 & 600 \\
\bottomrule
\end{tabular}
\end{table}

On the PD-only strict template benchmark, routing separation is
solved by the question templates. Therefore, any remaining
answer-quality gap comes from \emph{retrieval and evidence
construction} inside each route. The graph/hybrid handlers are
limited by low retrieval Hit@1 for canonical targets, so the LLM
often receives incomplete or missing evidence and answers become
paraphrastically shorter or partially ungrounded.

To contextualize the Stage~1 XGBoost classifier, we compare
it against several routing baselines trained and evaluated
on the same strict test split.
Table~\ref{tab:routing_baselines} reports the results.
The MajorityRoute baseline always predicts the most frequent
training class and achieves Macro-F1 $=0.222$, confirming
that simple frequency cannot solve routing. The
KeywordRuleRouter uses hand-crafted regex features and
reaches Macro-F1 $=0.497$. Logistic Regression and Random
Forest, trained on the same $28$-dimensional feature vector
as XGBoost, reach Macro-F1 $0.781$ and $0.790$ respectively.
The XGBoost row uses the same saved offline Stage~1 checkpoint
summarized in Table~\ref{tab:per_class_stage1}, giving
accuracy $0.805$, Macro-F1 $0.780$, and weighted-F1 $0.800$.
The PhoBERT sequence classifier achieves the highest
offline routing Macro-F1 $=0.901$, outperforming XGBoost
($0.780$). However, XGBoost is chosen as the Stage~1
classifier for two reasons: (1)~routing latency is
approximately $8.9$~ms for XGBoost versus $53.3$~ms for
PhoBERT per query, which matters when Stage~1 is a
gateway for every incoming query; and (2)~XGBoost runs on
CPU without a GPU dependency, simplifying deployment.
Although Logistic Regression and Random Forest obtain
comparable or slightly higher offline Macro-F1, XGBoost is
retained as the implemented Stage~1 checkpoint because it
supports non-linear interactions among the hand-crafted
routing features and provides the confidence score $s_1$ used
by the Stage~2 trigger. Since the differences among these
feature-based models are small, the main deployment trade-off
is between the lightweight feature-based router and the
higher-accuracy but slower PhoBERT classifier.
PhoBERT as a Stage~1 replacement is a promising direction
for future work where higher routing accuracy can justify
the additional inference cost.

\begin{table}[!htbp]
\centering
\caption{Routing classifier comparison on the strict test split ($600$ queries). Non-oracle classifiers are evaluated from the user question only: feature-based models use derived query-level routing features, while PhoBERT uses the raw question text. MetadataRuleRouter is oracle-style and excluded from comparison. PhoBERT achieves the best Macro-F1, but XGBoost is retained for the deployed Stage~1 router because of latency and deployment simplicity.}
\label{tab:routing_baselines}
\scriptsize
\setlength{\tabcolsep}{2.4pt}
\begin{tabular}{lrrr}
\toprule
\textbf{Classifier} & \textbf{Acc.} & \textbf{Macro-F1} & \textbf{W-F1} \\
\midrule
MajorityRoute & 0.500 & 0.222 & 0.333 \\
KeywordRuleRouter & 0.612 & 0.497 & 0.568 \\
Logistic Regression & 0.825 & 0.781 & 0.806 \\
Random Forest & 0.833 & 0.790 & 0.815 \\
XGBoost (Stage~1) & 0.805 & 0.780 & 0.800 \\
PhoBERT-base-v2 & \textbf{0.913} & \textbf{0.901} & \textbf{0.910} \\
\bottomrule
\end{tabular}
\end{table}

\begin{figure}[!htbp]
\centering
\begin{tikzpicture}
\begin{axis}[
    ybar,
    width=0.95\columnwidth,
    height=5.2cm,
    ymin=0, ymax=1.05,
    ylabel={Macro-F1},
    symbolic x coords={Majority,Keyword,LR,RF,XGBoost,PhoBERT},
    xtick=data,
    xticklabel style={rotate=30,anchor=east,font=\scriptsize},
    bar width=10pt,
    nodes near coords,
    nodes near coords style={font=\tiny},
    tick label style={font=\scriptsize},
    label style={font=\scriptsize},
    enlarge x limits=0.12
]
\addplot coordinates {
    (Majority,0.222)
    (Keyword,0.497)
    (LR,0.781)
    (RF,0.790)
    (XGBoost,0.780)
    (PhoBERT,0.901)
};
\end{axis}
\end{tikzpicture}
\caption{Routing classifier Macro-F1 on the strict test split.
PhoBERT achieves the best offline routing accuracy but at
higher inference cost ($53.3$~ms vs.\ $8.9$~ms for XGBoost).}
\label{fig:routing_baselines}
\end{figure}

Latency also supports the routing motivation. Single-stage
Router is slightly faster than Pure Graph ($2{,}209.2$ ms vs.
$2{,}283.4$ ms) while producing higher F1. This result is
important because the project is not trying to always use
GraphRAG; it is trying to avoid graph traversal when dense
retrieval is sufficient, while still preserving graph and
hybrid routes for relation-heavy questions.
The $3.3\%$ latency reduction over Pure Graph, while achieving
approximately $16.7\%$ higher F1 than Pure Vector, supports
the main motivation of adaptive routing: simple lookup queries
should not pay the cost of graph traversal.

\subsection{Clarification Benchmark}
\label{subsec:clarify_results}

Table~\ref{tab:clarify_results} reports the original
pre-Phase-3 configuration on the $234$-query constructed
ambiguity benchmark. The most important finding is that
Stage~1 has clarify F1 $=0.000$. This is not an
implementation bug. Stage~1 is trained on the strict
three-class dataset with labels \texttt{dense\_retrieval},
\texttt{graph\_traversal}, and \texttt{hybrid\_reasoning};
it has no \texttt{clarify} class during training and
therefore cannot architecturally predict clarification. In
this initial configuration, clarification behavior depends on
Stage~2 and the ambiguity trigger.

\begin{table}[!htbp]
\centering
\caption{Original pre-Phase-3 clarification results on the constructed ambiguity benchmark.}
\label{tab:clarify_results}
\scriptsize
\setlength{\tabcolsep}{2.2pt}
\resizebox{\columnwidth}{!}{%
\begin{tabular}{lrrrrr}
\toprule
\textbf{Variant} & \textbf{Route Acc.} & \textbf{P} & \textbf{R} & \textbf{F1} & \textbf{Trigger} \\
\midrule
Stage~1 only & 0.269 & 0.000 & 0.000 & 0.000 & 0.000 \\
Stage~1 + Stage~2 & 0.585 & 1.000 & 0.519 & 0.684 & 0.577 \\
\bottomrule
\end{tabular}%
}
\end{table}

The initial two-stage configuration improves clarify F1 from
$0.000$ to $0.684$. Its clarify precision is $1.000$, meaning
that when the system predicts \texttt{clarify}, the prediction
is correct in this benchmark. However, clarify recall is only
$0.519$, so many ambiguous queries are still missed. This
pattern is preferable to an over-clarifying system in a first
iteration because it avoids false clarification, but it also
shows that the initial ambiguity detector is incomplete.

Table~\ref{tab:ambiguity_type_results} explains the failure
mode. Stage~2 detects surface-context ambiguity very well:
unresolved pronoun references and incomplete context both
have $100\%$ trigger rate and $100\%$ recall. However, it
completely misses missing-entity ambiguity and multiple-
interpretation ambiguity in the current benchmark. These two
categories require more semantic understanding because the
query can look grammatically complete while still lacking the
specific legal target or allowing multiple valid legal
interpretations.

\begin{table}[!htbp]
\centering
\caption{Stage~2 behavior by ambiguity type.}
\label{tab:ambiguity_type_results}
\small
\setlength{\tabcolsep}{3.0pt}
\begin{tabular}{lrrr}
\toprule
\textbf{Ambiguity type} & \textbf{Total} & \textbf{Trigger} & \textbf{Recall} \\
\midrule
Pronoun reference & 42 & 1.000 & 1.000 \\
Incomplete context & 39 & 1.000 & 1.000 \\
Missing entity & 39 & 0.000 & 0.000 \\
Multi-interpretation & 36 & 0.000 & 0.000 \\
\bottomrule
\end{tabular}
\end{table}

\begin{figure}[!htbp]
\centering
\begin{tikzpicture}
\begin{axis}[
    ybar,
    width=0.95\columnwidth,
    height=5.0cm,
    ymin=0, ymax=1.15,
    ylabel={Score},
    symbolic x coords={Acc, Precision, Recall, F1},
    xtick=data,
    bar width=7pt,
    legend style={at={(0.5,1.05)},anchor=south,
                  legend columns=2,font=\scriptsize},
    tick label style={font=\scriptsize},
    label style={font=\scriptsize},
    nodes near coords,
    nodes near coords style={font=\tiny,rotate=90,anchor=west},
    enlarge x limits=0.2
]
\addplot coordinates {(Acc,0.269)(Precision,0.000)
                      (Recall,0.000)(F1,0.000)};
\addplot coordinates {(Acc,0.585)(Precision,1.000)
                      (Recall,0.519)(F1,0.684)};
\legend{Stage~1 only, Stage~1+Stage~2}
\end{axis}
\end{tikzpicture}
\caption{Original pre-Phase-3 clarification metrics on the ambiguity benchmark.
Stage~1 scores zero on all clarification metrics because it has
no clarify label. Stage~2 achieves perfect precision ($1.000$)
by relying on explicit surface-context cues.}
\label{fig:clarify_f1_bar}
\end{figure}

\subsection{Conversation-aware Ambiguity Stress Test}
\label{subsec:conversation_ambiguity_results}

After adding deterministic history-referent resolution and a
severe-ambiguity override policy, we evaluate the router on the
$160$-query conversation stress test described in
Section~\ref{subsec:dataset_description}. This diagnostic is
reported separately from the strict end-to-end benchmark because
it measures whether the router uses conversation history
correctly, not answer generation quality.
Table~\ref{tab:conversation_stress_summary} compares the router
before and after Phase~3. Clarify F1 improves from $0.545$ to
$0.862$, and history-resolution accuracy reaches $0.850$. The
improvement comes from two complementary behaviors: resolved
history can unblock retrieval, while absent, irrelevant, or
conflicting history can preserve or force clarification.

\begin{table}[!htbp]
\centering
\caption{Diagnostic conversation-aware ambiguity stress-test results.}
\label{tab:conversation_stress_summary}
\scriptsize
\setlength{\tabcolsep}{2.2pt}
\resizebox{\columnwidth}{!}{%
\begin{tabular}{lrrrrr}
\toprule
\textbf{Variant} & \textbf{Route Acc.} & \textbf{P} & \textbf{R} & \textbf{F1} & \textbf{Hist. Acc.} \\
\midrule
Before Phase~3 & 0.475 & 0.907 & 0.390 & 0.545 & n/a \\
After Phase~3 & 0.750 & 0.963 & 0.780 & 0.862 & 0.850 \\
\bottomrule
\end{tabular}%
}
\end{table}

Table~\ref{tab:conversation_category_results} reports the
after-Phase~3 behavior by category. Clear dense controls
remain intact with route accuracy $1.000$ and a low Stage~2
rate, showing that the dense fast path is preserved. The
router also improves clarification for no-history, irrelevant
history, conflicting history, missing-entity, and
multi-interpretation cases. However, answerable follow-up
questions and graph/hybrid controls remain difficult because
the router may resolve that a query is answerable but still
confuse the correct retrieval mechanism.

\begin{table}[!htbp]
\centering
\caption{Diagnostic conversation stress-test category results after Phase~3.}
\label{tab:conversation_category_results}
\scriptsize
\setlength{\tabcolsep}{2.0pt}
\begin{tabularx}{\columnwidth}{L{2.45cm}rrY}
\toprule
\textbf{Category} & \textbf{Score} & \textbf{S2} & \textbf{Interpretation} \\
\midrule
Answerable history & 0.600 & 0.850 & Remaining dense/graph/hybrid confusion. \\
No-history clarify & 0.850 & 1.000 & No-history references mostly clarified. \\
Irrelevant history & 0.750 & 0.900 & Irrelevant history no longer suppresses ambiguity. \\
Conflicting history & 0.700 & 0.900 & Multiple candidates remain difficult. \\
Missing entity & 0.850 & 1.000 & Missing legal targets are better detected. \\
Multi-interpret. & 0.750 & 0.850 & Improved on the conversation stress test. \\
Dense control & 1.000 & 0.150 & Dense fast path is preserved. \\
Graph/hybrid ctrl. & 0.500 & 0.950 & Relation-heavy route choice remains hard. \\
\bottomrule
\end{tabularx}
\end{table}

The original $234$-query ambiguity benchmark was also rerun
after adding deterministic history-referent resolution and the
severe-ambiguity override, as shown in
Table~\ref{tab:clarify_phase3_before_after}. Clarify F1
increases from $0.684$ to $0.840$ with zero false positives.
The strongest gains are on missing-entity cases, but
multi-interpretation remains weak on the original template
benchmark because these queries often look answerable without
an explicit contextual pronoun.

\begin{table}[!htbp]
\centering
\caption{Original ambiguity benchmark before and after deterministic history-referent resolution and severe-ambiguity override.}
\label{tab:clarify_phase3_before_after}
\scriptsize
\setlength{\tabcolsep}{2.5pt}
\resizebox{\columnwidth}{!}{%
\begin{tabular}{lrrrrr}
\toprule
\textbf{Variant} & \textbf{Acc.} & \textbf{P} & \textbf{R} & \textbf{F1} & \textbf{FP} \\
\midrule
Before Phase~3 & 0.585 & 1.000 & 0.519 & 0.684 & 0 \\
After Phase~3 & 0.722 & 1.000 & 0.724 & 0.840 & 0 \\
\bottomrule
\end{tabular}%
}
\end{table}

Qualitatively, the demo shows the intended behavior. With
history mentioning ``\textit{Nghị định 100/2019/NĐ-CP}'', the
follow-up query ``\textit{Văn bản đó còn hiệu lực không?}''
is resolved to that decree and routed to graph traversal.
With empty or irrelevant history, the same surface query is
routed to \texttt{clarify} instead of retrieving evidence for
an arbitrary document.
Table~\ref{tab:qualitative_routing_cases} provides four
concrete examples illustrating the intended routing behavior
on these cases.

\begin{table}[!htbp]
\centering
\caption{Qualitative routing examples illustrating the intended
behavior of the adaptive router. These examples are diagnostic
and are not counted as additional benchmark results.}
\label{tab:qualitative_routing_cases}
\scriptsize
\setlength{\tabcolsep}{2.0pt}
\renewcommand{\arraystretch}{1.08}
\begin{tabularx}{\columnwidth}{L{2.15cm}L{2.70cm}Y}
\toprule
\textbf{Case} & \textbf{Example query / context} & \textbf{Router behavior} \\
\midrule
Direct lookup &
\textit{Điều kiện cấp giấy phép xây dựng là gì?} &
The query asks for legal conditions stated in text, so the
router can select dense retrieval without paying graph
traversal cost. \\
Relation-heavy &
\textit{Văn bản nào sửa đổi Nghị định 100/2019?} &
The amendment cue indicates a legal relation, so graph
traversal or hybrid reasoning is more appropriate than pure
dense lookup. \\
Unresolved reference &
\textit{Văn bản đó còn hiệu lực không?}, with $H=\emptyset$ &
The referent is missing. The safer action is clarification
rather than retrieving evidence for an arbitrary document. \\
Resolved history &
History mentions \textit{Nghị định 100/2019/NĐ-CP}; current
query: \textit{Văn bản đó còn hiệu lực không?} &
The history resolver can bind \textit{văn bản đó} to a unique
legal document, allowing the router to proceed with retrieval
instead of unnecessary clarification. \\
\bottomrule
\end{tabularx}
\end{table}

\subsection{Stage~2 Cost and Override Behavior}
\label{subsec:stage2_cost}

The benefit of Stage~2 comes with a substantial latency cost.
In the strict full benchmark, Two-stage Hybrid triggers
Stage~2 for $43.0\%$ of queries and changes the Stage~1 route
for only $4.26\%$ of triggered cases. In the ambiguity
benchmark, Stage~2 is triggered for $57.7\%$ of queries and
has a much higher override rate of $69.6\%$. These two
numbers are not contradictory. The strict benchmark contains
mostly answerable retrieval questions, whereas the ambiguity
benchmark intentionally contains many underspecified queries.

Routing-only latency measurements show the same pattern.
Stage~1-only routing costs approximately $33.7$ ms on
$342$ non-Stage~2 cases, while routing with Stage~2 costs
approximately $4{,}017$ ms on $258$ Stage~2 cases. This is a
roughly $119\times$ latency increase. Therefore, invoking
Stage~2 for every query would undermine the efficiency goal
of adaptive routing. The correct design is selective
verification: Stage~1 should handle clear cases, while
Stage~2 should be reserved for ambiguous, uncertain, or
relation-heavy queries.

The Phase~3 conversation stress test confirms this trade-off.
Clarification quality improves, but Stage~2 is triggered for
$82.5\%$ of conversation-stress queries and average routing
latency is $3{,}127.2$ ms. A separate routing-only strict
sanity run after Phase~3 reaches routing accuracy $0.8933$
with $10$ false clarifications out of $600$ queries. Because
this sanity run does not perform retrieval or answer
generation, and because it is lower than the full strict
end-to-end routing result in Table~\ref{tab:end_to_end_results},
it is not used to replace the main strict benchmark results.

\begin{figure}[!htbp]
\centering
\begin{tikzpicture}
\begin{axis}[
    ybar,
    width=0.95\columnwidth,
    height=4.8cm,
    ymin=0,
    ylabel={Avg.\ latency (ms)},
    symbolic x coords={PureVec,PureGraph,SingleStg,TwoStg},
    xtick=data,
    xticklabels={{Pure\\Vector},{Pure\\Graph},
                 {Single-\\stage},{Two-\\stage}},
    xticklabel style={align=center,font=\scriptsize},
    bar width=14pt,
    nodes near coords,
    nodes near coords style={font=\scriptsize},
    tick label style={font=\scriptsize},
    label style={font=\scriptsize},
    enlarge x limits=0.18
]
\addplot coordinates {
    (PureVec,1270.7)
    (PureGraph,2283.4)
    (SingleStg,2209.2)
    (TwoStg,3913.4)
};
\end{axis}
\end{tikzpicture}
\caption{Average end-to-end latency per query. Single-stage Router
is faster than Pure Graph while achieving higher F1, confirming
that adaptive routing reduces unnecessary graph traversal. Two-stage
Hybrid latency reflects the $43\%$ Stage~2 trigger rate.}
\label{fig:latency_comparison}
\end{figure}

\subsection{Threshold Ablation}
\label{subsec:ablation_results}

Table~\ref{tab:threshold_ablation} summarizes the threshold
ablation on the development setting. Sweeping the Stage~1
confidence threshold $\tau_{\mathrm{conf}}$ from $0.50$ to $0.99$ does not
change routing accuracy, which remains $0.7529$. It only
changes Stage~2 trigger rate. This indicates that many
Stage~1 predictions are already high confidence, so
$\tau_{\mathrm{conf}}$ is not the main bottleneck in this setting.

\begin{table}[!htbp]
\centering
\caption{Selected threshold ablation results.}
\label{tab:threshold_ablation}
\small
\setlength{\tabcolsep}{3.5pt}
\begin{tabular}{lrrr}
\toprule
\textbf{Parameter} & \textbf{Value} & \textbf{Acc.} & \textbf{Stage2 rate} \\
\midrule
$\tau_{\mathrm{conf}}$ & 0.50 & 0.7529 & 0.2299 \\
$\tau_{\mathrm{conf}}$ & 0.70 & 0.7529 & 0.3046 \\
$\tau_{\mathrm{conf}}$ & 0.95 & 0.7529 & 0.5920 \\
$\tau_{\mathrm{conf}}$ & 0.99 & 0.7529 & 0.5977 \\
\midrule
$\tau_{\mathrm{clar}}$ & 0.50 & 0.6839 & 0.2299 \\
$\tau_{\mathrm{clar}}$ & 0.70 & 0.7529 & 0.2299 \\
$\tau_{\mathrm{clar}}$ & 0.90 & 0.7759 & 0.2299 \\
$\tau_{\mathrm{clar}}$ & 0.95 & 0.8448 & 0.2299 \\
\bottomrule
\end{tabular}
\end{table}

The clarification override threshold $\tau_{\mathrm{clar}}$ has a stronger effect. A
low value such as $0.50$ is too aggressive on non-ambiguous
development data and reduces routing accuracy to $0.6839$.
Increasing $\tau_{\mathrm{clar}}$ to $0.95$ improves routing accuracy to
$0.8448$. This result should be interpreted carefully: the
ablation is performed on a routing-oriented development set,
not on the ambiguity benchmark. Therefore, the threshold that
maximizes strict routing accuracy may not maximize clarify
recall. The practical conclusion is that $\tau_{\mathrm{clar}}$ must be
calibrated jointly for both false clarification on answerable
queries and missed clarification on ambiguous queries.

\subsection{Overall Discussion}
\label{subsec:overall_discussion}

A useful way to interpret the results is that the proposed
router is not intended to make graph retrieval stronger in
isolation. Its role is to decide when graph retrieval is
actually warranted. Graph evidence can improve relational
reasoning, but it can also increase cost or introduce broad
context when used indiscriminately. The strict benchmark
therefore supports a selective design: Pure Graph is not
better than Pure Vector overall, but a router that selects
among dense, graph, and hybrid routes improves answer F1 and
routing accuracy. The additional \texttt{clarify} route
extends this selection view to legal safety, where the best
decision for an underspecified query may be to postpone
retrieval and ask for the missing legal target.

The experiments support the main thesis of the project:
Vietnamese legal QA should not use a fixed retrieval strategy
for every query. Dense retrieval is strong for direct lookup,
but it is weak when the answer depends on explicit legal
relations. Graph traversal is useful for relation-heavy
queries, but applying it to every question can add cost and
produce diffuse context. Hybrid reasoning is appropriate when
both relation paths and article text are needed. Clarification
is necessary when the legal target is underspecified.

The current strongest result is not that Stage~2 always
improves every metric. Instead, the stronger and more honest
finding is that Stage~1 is already sufficient for many strict
retrieval-routing cases, while Stage~2 is necessary for
clarification because Stage~1 has no clarify label. The
clarify benchmark reveals a precise limitation: Stage~2
recognizes surface contextual ambiguity but misses semantic
ambiguity. This finding is useful for the next development
step because it points directly to missing-entity detection,
entity-linking uncertainty, and multi-interpretation
generation as the components that should be improved.

The conversation-aware stress test strengthens this analysis.
The router no longer treats conversation history as a binary
feature. Instead, it distinguishes resolved history from
absent, irrelevant, and conflicting history. This allows a
follow-up query with a unique referent to proceed to retrieval,
while the same query without a usable referent can ask for
clarification. The trade-off is higher Stage~2 usage and
latency on conversation-stress examples. Moreover, the
routing-only strict sanity run after Phase~3 still produces
$10$ false clarifications, so the Phase~3 diagnostic results
should be interpreted as an ambiguity improvement rather than
a replacement for the strict end-to-end QA table.

% ============================================================
\section{Limitations and Future Work}
\label{sec:limitations}
% ============================================================

Several limitations remain. First, the strict training split
does not include a \texttt{clarify} class, so Stage~1 cannot
predict clarification by design. The initial Stage~2
configuration improves clarify F1 from $0.000$ to $0.684$,
and the Phase~3 additions improve the original ambiguity
benchmark further to $0.840$ and the separate conversation
stress test to $0.862$. However, clarification still depends
on the LLM verifier and remains sensitive to trigger
calibration. Second, the ambiguity benchmarks are synthetic
and template-generated. They are useful for controlled stress
testing, but they should be validated with manually written or
human-reviewed ambiguous legal questions before being treated
as realistic production benchmarks. Third, the initial
configuration detects surface cues such as unresolved pronouns
and incomplete context more reliably than semantic ambiguity.
After Phase~3, missing-entity cases improve substantially, but
multi-interpretation remains weak on the original template
benchmark.

Fourth, Stage~2 introduces a large latency cost. Routing with
Stage~2 is about $119\times$ slower than Stage~1-only routing
in the measured routing logs. Therefore, the verifier should
not be used for every query. Future work should calibrate the
trigger function with both strict routing and ambiguity
benchmarks, rather than optimizing only one setting. Fifth,
PD-only strict routing accuracy remains at $1.0000$ even when Stage~2
is enabled. Stage~2 therefore does not fix route choice on this split;
it mainly adds latency, while the remaining end-to-end gap is driven by
weaker graph/hybrid evidence retrieval.

The Phase~3 conversation-aware extension introduces additional
limitations. Although conversation clarify F1 improves to
$0.862$, \texttt{answerable\_with\_history} route accuracy is
only $0.600$, meaning that the system can resolve that a query
is answerable but still confuse dense, graph, and hybrid
execution. The clear graph/hybrid control category also
remains difficult with route accuracy $0.500$. On the original
template ambiguity benchmark, \texttt{multi\_interpretation}
remains weak even after Phase~3, despite improvements on
missing-entity and pronoun-reference cases. Finally, a
routing-only strict sanity run after Phase~3 still produces
$10$ false clarifications out of $600$ queries; therefore,
these diagnostic improvements are not used to replace the
strict end-to-end metrics in Table~\ref{tab:end_to_end_results}.

Sixth, the graph is larger and better connected after the
migration, but curated legal relations remain sparse:
\textsc{AMENDS}, \textsc{REPEALS}, and \textsc{GUIDES} have
far fewer edges than generic \textsc{REFERENCES}. As a
result, graph traversal may still miss important legal
relations or return overly broad reference context. Seventh,
the primary routing labels are derived from hop-count and
cross-document metadata rather than independent human
annotation. These labels are useful for training and
benchmarking, but manual review is needed to validate whether
the assigned route truly matches the best retrieval strategy.
Finally, end-to-end answer quality is evaluated with
automatic string-based metrics, which are imperfect for legal
answers. Future evaluation should include human judgments of
legal correctness, evidence faithfulness, citation usefulness,
and whether clarification questions are appropriate.

Future work will focus on integrating manually reviewed
\texttt{clarify} samples into the main training set,
improving semantic ambiguity detection through entity-linking
uncertainty and candidate-interpretation generation,
enriching the legal graph with more reliable amendment and
repeal relations, investigating a PhoBERT-based Stage~1
router to improve offline routing Macro-F1 from $0.780$ to
approximately $0.901$ when the additional inference latency
is acceptable, applying class-balancing techniques to address
the training-set imbalance toward \texttt{dense\_retrieval},
and expanding answer-quality evaluation with human review.

% ============================================================
\section{Conclusion}
\label{sec:conclusion}
% ============================================================

This paper presented a Reasoning-aware Adaptive Router for
Vietnamese Legal Hybrid GraphRAG. The system combines a dense
vector index, a Neo4j legal knowledge graph, a Stage~1 XGBoost
router, and a selective Stage~2 LLM verifier. Instead of
always using dense retrieval or always using graph traversal,
the router selects among dense retrieval, graph traversal,
hybrid reasoning, and clarification according to the query's
reasoning demand and ambiguity level.

On the PD-only strict $600$-query benchmark, the adaptive router
selects the correct route (routing accuracy $1.0000$ for both
single-stage and two-stage systems), but end-to-end F1 is still
dominated by dense retrieval because graph/hybrid evidence
retrieval quality is lower. Pure Vector reaches F1 $=0.4947$
while Pure Graph reaches $0.3602$; Single-stage Router and
Two-stage Hybrid are close to Pure Graph (F1 $\approx 0.3610/0.3602$).
This indicates that Stage~2 should be interpreted primarily as
a selective verifier/ambiguity safety layer, not as a mechanism
that guarantees F1 improvements when the route's retrieved
evidence is weak.

The ambiguity benchmarks show that clarification handling is
the main role of Stage~2. Stage~1 has clarify F1 $=0.000$
because it is trained only on three retrieval labels. The
initial Stage~2 configuration improves clarify F1 to $0.684$
with precision $1.000$. After adding deterministic
history-referent resolution and the severe-ambiguity override,
the original ambiguity benchmark improves to F1 $=0.840$.
The Phase~3 conversation-aware extension further shows that
history resolution is a useful intermediate signal for legal
QA: on the separate $160$-query conversation stress test, the
router reaches clarify F1 $=0.862$ and history-resolution
accuracy $=0.850$.

The same follow-up query can therefore be routed differently
depending on whether history contains a unique legal referent.
However, multi-interpretation ambiguity remains weak on the
original template benchmark, resolved-history route selection
is still imperfect, and the strict routing-only sanity check
still produces false clarifications. These results support the
central design of reasoning-aware adaptive routing while
showing that the Phase~3 diagnostic improvements should not
replace the strict end-to-end configuration without further
calibration.

\section*{Acknowledgements}

The author thanks the School of Information Technology and
Communication at Hanoi University of Science and Technology
for supporting this research. The author is grateful to
supervisors Pham Van Hai and Nguyen Viet Cuong for their
guidance throughout the project.

\bibliographystyle{spmpsci}
\bibliography{biblio}

\end{document}
